\ifdefined\gkver\else\def\gkver{student}\fi
\documentclass[\gkver]{gaokaozhenti}
\begin{document}

\chapter{2005年辽宁文理综合}
%% paperType: 文理综合物理部分

\begin{enumerate}
\item
%% number: 31
%% paperName: 2005年辽宁文理综合第 31 题 6 分
%% typeId: 1
%% type: 单选题
%% chapter: 原子和原子核
%% point: 天然放射性
%% method: 无
%% score: 6
%% degree: 750
%% duplicateId: 0
%% body:
如图所示，放射源放在铅块上的细孔中，铅块上方有匀强磁场，磁场方向垂直于纸面向外。已知放射源放出的射线有 $\alpha$、$\beta$、$\gamma$ 三种。下列判断正确的是\xzanswer{B}

\onepicture[w=3.4cm]{figs/31.png}

\fourchoices[answer=B]
{甲是 $\alpha$ 射线，乙是 $\gamma$ 射线，丙是 $\beta$ 射线}
{甲是 $\beta$ 射线，乙是 $\gamma$ 射线，丙是 $\alpha$ 射线}
{甲是 $\gamma$ 射线，乙是 $\alpha$ 射线，丙是 $\beta$ 射线}
{甲是 $\alpha$ 射线，乙是 $\beta$ 射线，丙是 $\gamma$ 射线}

%% answer: B
%% memo:
\memoanswer{
本题原卷及材料中未提供详解，待补充。
}

\item
%% number: 32
%% paperName: 2005年辽宁文理综合第 32 题 6 分
%% typeId: 1
%% type: 单选题
%% chapter: 电路
%% point: 多用表的使用
%% method: 无
%% score: 6
%% degree: 750
%% duplicateId: 0
%% body:
图中 B 为电源，$R_{1}$、$R_{2}$ 为电阻，K 为电键。现用多用电表测量流过电阻 $R_{2}$ 的电流。将多用电表的选择开关调至直流电流挡（内阻很小）以后，正确的接法是\xzanswer{C}

\onepicture[w=4.4cm]{figs/32.png}

\fourchoices[answer=C]
{保持 K 闭合，将红表笔接在 a 处，黑表笔接在 b 处}
{保持 K 闭合，将红表笔接在 b 处，黑表笔接在 a 处}
{将 K 断开，红表笔接在 a 处，黑表笔接在 b 处}
{将 K 断开，红表笔接在 b 处，黑表笔接在 a 处}

%% answer: C
%% memo:
\memoanswer{
本题原卷及材料中未提供详解，待补充。
}

\item
%% number: 33
%% paperName: 2005年辽宁文理综合第 33 题 6 分
%% typeId: 1
%% type: 单选题
%% chapter: 光学
%% point: 光的折射
%% method: 无
%% score: 6
%% degree: 650
%% duplicateId: 0
%% body:
一束复色光由空气射向一块平行平面玻璃砖，经折射分成两束单色光 a、b。已知 a 光的频率小于 b 光的频率。下列哪个光路图可能是正确的？\xzanswer{B}

\fourchoices[answer=B, ispicture=true, h=1.8cm]
{figs/33a.png}
{figs/33b.png}
{figs/33c.png}
{figs/33d.png}

%% answer: B
%% memo:
\memoanswer{
本题原卷及材料中未提供详解，待补充。
}

\item
%% number: 34
%% paperName: 2005年辽宁文理综合第 34 题 6 分
%% typeId: 1
%% type: 单选题
%% chapter: 电磁感应
%% point: 电磁感应电路分析
%% method: 无
%% score: 6
%% degree: 650
%% duplicateId: 0
%% body:
如图所示，两根相距为 $l$ 的平行直导轨 ab、cd，b、d 间连有一固定电阻 $R$，导轨电阻可忽略不计。MN 为放在 ab 和 cd 上的一导体杆，与 ab 垂直，其电阻也为 $R$。整个装置处于匀强磁场中，磁感应强度的大小为 $B$，磁场方向垂直于导轨所在平面（指向图中纸面内）。现对 MN 施力使它沿导轨方向以速度 $v$（如图）做匀速运动。令 $U$ 表示 MN 两端电压的大小，则\xzanswer{A}

\onepicture[w=6.2cm]{figs/34.png}

\fourchoices[answer=A]
{$U=\frac{1}{2}Blv$，流过固定电阻 $R$ 的感应电流由 b 到 d}
{$U=\frac{1}{2}Blv$，流过固定电阻 $R$ 的感应电流由 d 到 b}
{$U=Blv$，流过固定电阻 $R$ 的感应电流由 b 到 d}
{$U=Blv$，流过固定电阻 $R$ 的感应电流由 d 到 b}

%% answer: A
%% memo:
\memoanswer{
本题原卷及材料中未提供详解，待补充。
}

\item
%% number: 35
%% paperName: 2005年辽宁文理综合第 35 题 6 分
%% typeId: 1
%% type: 单选题
%% chapter: 功和能
%% point: 动能定理
%% method: 无
%% score: 6
%% degree: 650
%% duplicateId: 0
%% body:
一物块由静止开始从粗糙斜面上的某点加速下滑到另一点，在此过程中重力对物块做的功等于\xzanswer{D}

\fourchoices[answer=D]
{物块动能的增加量}
{物块重力势能的减少量与物块克服摩擦力做的功之和}
{物块重力势能的减少量和物块动能的增加量以及物块克服摩擦力做的功之和}
{物块动能的增加量与物块克服摩擦力做的功之和}

%% answer: D
%% memo:
\memoanswer{
本题原卷及材料中未提供详解，待补充。
}

\item
%% number: 36
%% paperName: 2005年辽宁文理综合第 36 题 6 分
%% typeId: 1
%% type: 单选题
%% chapter: 力物体的平衡
%% point: 三力平衡
%% method: 无
%% score: 6
%% degree: 450
%% duplicateId: 0
%% body:
两光滑平板 MO、NO 构成一具有固定夹角 $\theta_{0}=75^{\circ}$ 的 V 形槽，一球置于槽内，用 $\theta$ 表示 NO 板与水平面之间的夹角，如图所示。若球对板 NO 压力的大小正好等于球所受重力的大小，则下列 $\theta$ 值中哪个是正确的？\xzanswer{B}

\onepicture[w=6.6cm]{figs/36.png}

\fourchoices[answer=B]
{$15^{\circ}$}
{$30^{\circ}$}
{$45^{\circ}$}
{$60^{\circ}$}

%% answer: B
%% memo:
\memoanswer{
本题原卷及材料中未提供详解，待补充。
}

\item
%% number: 39
%% paperName: 2005年辽宁文理综合第 39 题 14 分
%% typeId: 6
%% type: 计算题
%% chapter: 电场
%% point: 电势能电势差电场力做功
%% method: 无
%% score: 14
%% degree: 550
%% duplicateId: 0
%% body:
一匀强电场，场强方向是水平的（如图）。一个质量为 $m$ 的带正电的小球，从 O 点出发，初速度的大小为 $v_{0}$，在电场力与重力的作用下，恰能沿与场强的反方向成 $\theta$ 角的直线运动。求小球运动到最高点时其电势能与在 O 点的电势能之差。

\onepicture[w=7.6cm]{figs/39.png}

%% answer: 
\jdanswer{
$\Delta E_{p}=\frac{1}{2}mv_{0}^{2}\cos^{2}\theta$
}
%% memo:
\memoanswer{
设电场强度为 $E$，小球带电量为 $q$，因小球做直线运动，它受的电场力 $qE$ 和重力 $mg$ 的合力必沿此直线，如图。

$\frac{mg}{qE}=\tan\theta$。（4 分）

由此可知，小球做匀减速运动的加速度大小为

$a=\frac{g}{\sin\theta}$。（2 分）

设从 O 到最高点的路程为 $s$，

$v_{0}^{2}=2as$。（2 分）

运动的水平距离为 $l=s\cos\theta$。（2 分）

两点的电势能之差 $\Delta E_{p}=qEl$。（2 分）

\onepicture[w=7.6cm]{figs/39_an.png}

由以上各式得

$\Delta E_{p}=\frac{1}{2}mv_{0}^{2}\cos^{2}\theta$。（2 分）
}

\end{enumerate}

\end{document}
